% ============================================================================
% ButterflAI: A hybrid statistical + score-based generative model of the
% solar butterfly diagram.
%
% FIRST DRAFT. Prose written to the author's voice profile (see voice/).
% Citations are placeholders [CIT] keyed to deep_research_text_markdown.md;
% the author supplies the final .bib.  Passages marked % TODO(andres) require
% the author's stance rather than a fabricated one.
%
% Primary class is AASTeX; an article-class fallback is provided in a comment
% block below in case AASTeX is not installed locally.
% ============================================================================
\documentclass[twocolumn]{aastex631}

% ---- article-class fallback (uncomment if aastex631 is unavailable) --------
% \documentclass[11pt]{article}
% \usepackage{amsmath,amssymb,graphicx,natbib}
% \usepackage[margin=1in]{geometry}
% \newcommand{\facility}[1]{}
% ---------------------------------------------------------------------------

\usepackage{amsmath}

\newcommand{\degr}{\ensuremath{^{\circ}}}
\newcommand{\msh}{\ensuremath{\mu\mathrm{Hem}}}   % millionths of a solar hemisphere
\newcommand{\tzero}{\ensuremath{t_0}}
\newcommand{\mupeak}{\ensuremath{\mu_{\mathrm{peak}}}}
\newcommand{\muzero}{\ensuremath{\mu_0}}

\shorttitle{ButterflAI}
\shortauthors{Mu\~noz-Jaramillo et al.}

\begin{document}

\title{ButterflAI: A Hybrid Statistical and Score-Based Generative Model
of the Solar Butterfly Diagram}

% TODO(andres): author list and affiliations.
\author{A.~Mu\~noz-Jaramillo}
\affiliation{Southwest Research Institute, Boulder, CO, USA}

\author{The ButterflAI Team}
\affiliation{COFFIES Science Center}

\begin{abstract}
The equatorward migration of the sunspot emergence zones (the solar butterfly
diagram) is one of the most recognizable signatures of the solar cycle, yet
we lack a compact generative model able to draw new, physically plausible
butterfly diagrams from a small number of cycle parameters. We build such a
model in two parts. First, we describe a single butterfly wing as a
time-indexed family of Gaussians in absolute latitude, with a mean latitude
that follows a universal exponential path $\mu(\tau)=a\,e^{-\tau/b}$ (with
$a=16.31\degr$ and $b=8.14$~yr) once time is measured from a per-hemicycle
reference epoch, and with a width $\sigma(\mu)$ whose equatorward arm collapses
onto a single line shared by every hemicycle while its poleward arm scales with
cycle amplitude. Leave-one-cycle-out validation places the generalization floor
of this classical model at $3.087$ nats per year-block. Second, we train a
score-based diffusion model to generate the residual density that the classical
Gaussian leaves behind, conditioned on cycle amplitude and phase. On six
hemispheric cycles held out from all fitting, the hybrid model improves
distributional fidelity substantially over both the classical model and an
independent SARG-style baseline (a $24\%$ lower Earth Mover's distance and a
roughly halved spread error), although it does not improve the pointwise
likelihood, which continues to favor the smoother parametric Gaussian. We
interpret this tension as evidence that likelihood alone is the wrong sole
arbiter for a generative model of the butterfly diagram, and we show, using an
oracle diagnostic, that the residual signal available in the present
conditioning set is nearly exhausted. To the best of our knowledge, this is the
first application of a score-based generative model to the solar butterfly
diagram.
\end{abstract}

\keywords{Sunspots --- Solar cycle --- Solar activity --- Generative models}

% ============================================================================
\section{Introduction}
\label{sec:intro}
% ============================================================================

The butterfly diagram is the latitude-time distribution of sunspot emergence,
and its shape encodes much of what we know about the solar cycle: the
equatorward drift of the active latitudes (Sp\"orer's law), the $\sim$11-year
Schwabe periodicity, the amplitude relationships that connect a cycle's rise to
its strength, and the hemispheric asymmetries that make no two cycles identical
\citep[CIT]{}. A model that could \emph{generate} new butterfly diagrams, rather
than merely describe the observed ones, would let us build large synthetic
ensembles of solar cycles (thousands of realizations from a handful of
parameters), propagate cycle-to-cycle statistics into flux-transport and
irradiance reconstructions, and stress-test cycle-prediction methods against
populations of plausible-but-unobserved cycles. Such a
generator does not yet exist for the butterfly diagram itself.

Two distinct notions of a ``universal'' cycle motivate the model we build here,
and it is worth keeping them separate because they pull in opposite directions.
The first is the amplitude-\emph{independent} universality of the latitudinal
drift: \citet[CIT]{Hathaway2011} showed, from the area-weighted centroid of the
sunspot zones, that ``the sunspot zone latitudes and equatorward drift measured
relative to this starting time follow a standard path for all cycles with no
dependence upon cycle strength or hemispheric dominance,'' a result reinforced
by the universal decline phase of \citet[CIT]{CameronSchussler2016} and by the
mean butterfly diagram of \citet[CIT]{CameronSchussler2024}. The second is the
amplitude-\emph{dependent} universality of cycle \emph{shape}, expressed through
the Waldmeier effect and the one- or two-parameter shape functions that collapse
the sunspot-number time series onto a common curve
\citep[CIT]{HathawayWilsonReichmann1994, Volobuev2009}. However, there is a
normalization subtlety that bridges the two and which we exploit below: the
drift is universal only when time is measured from cycle \emph{onset}, and
systematic amplitude-dependent differences reappear when it is measured from
cycle \emph{maximum} \citep[CIT]{Hathaway2011}.

The nearest existing precedent for a synthetic butterfly generator is the
Synthetic Active Region Generator \citep[SARG;][CIT]{JhaUpton2024, Jha2025},
which draws a catalog of active regions (each with an emergence time, latitude,
flux, and tilt) from observed statistical relationships and feeds them to an
advective flux-transport model. SARG is a powerful tool, but it differs from
what we set out to build in three ways: it generates a discrete active-region
catalog rather than a direct latitude-time density; it fixes the emergence
centroid at $28\degr$ and ties only the \emph{width} of the emergence band to a
single, whole-Sun sunspot number; and it samples from fixed relationships
without a learned refinement of what those relationships leave behind. The
semi-synthetic reconstruction of \citet[CIT]{Jiang2011}, in contrast, does let
the mean emergence latitude rise with cycle amplitude ($\langle\lambda\rangle =
12.2\degr + 0.022\,S_n$), which is closer in spirit to the per-hemisphere
amplitude tying we adopt. Machine learning has touched the butterfly diagram
only as a forecasting tool so far \citep[CIT]{Covas2017, Covas2019}, and, to the
best of our knowledge, no diffusion, score-based, or other deep generative model
has been applied to the butterfly diagram itself, even though such models are
now well established elsewhere in solar physics for full-disk imagery and
active-region patches \citep[CIT]{Ramunno2024}.

In this work we build a generative model of the butterfly diagram in two parts.
The first part is a compact classical description of a single wing, in which the
wing is a family of Gaussians whose mean follows a universal exponential path
and whose width is a piecewise-linear function of latitude, with only a handful
of numbers left free per cycle. The second part is a score-based diffusion model
that learns the residual between the empirical latitude density and this
classical Gaussian, conditioned on the same cycle parameters the classical model
uses. We treat each hemisphere as its own amplitude-tied process throughout,
which distinguishes our model from generators that symmetrize the two
hemispheres or drive them from a single shared index.

In Section~\ref{sec:data} we describe the composite sunspot-group catalog and
the cuts we apply. In Section~\ref{sec:classical} we build the classical
universal-wing model, fit it, and validate it by leave-one-cycle-out
cross-validation and bootstrap. In Section~\ref{sec:diffusion} we set up the
score-based residual model and the metrics we use to evaluate it. In
Section~\ref{sec:results} we present the ablation, the failure of classifier-free
guidance, the information ceiling we uncover, and the held-out test comparison
against the classical and SARG baselines. In Section~\ref{sec:discussion} we
discuss the tension between likelihood and distributional fidelity, what the
residual framing assumes, and where the model can still be improved. We
summarize in Section~\ref{sec:summary}.

% ============================================================================
\section{Data}
\label{sec:data}
% ============================================================================

Our analysis rests on a composite catalog of daily sunspot-group measurements
spanning 1825--2023 (roughly $3.1\times10^{5}$ individual measurements), stitched
from several heliographic-position surveys into a single record with columns for
date, heliographic latitude and longitude, corrected area (in millionths of a
solar hemisphere, \msh, with an associated uncertainty), and solar cycle number.
Because the catalog is a composite of surveys with different (and not always
specified) detection thresholds, its early portion carries a heterogeneous
calibration; we therefore restrict all modeling to cycles $\geq 12$, the
well-observed and comparatively homogeneous era, while retaining the earlier
cycles only for context.

We treat each hemisphere of each cycle as an independent unit, which we refer to
from now on as a \emph{hemicycle}, giving us $26$ hemicycles over cycles 12--24.
This choice doubles the effective sample and imposes the (approximate)
north-south symmetry of the emergence process, and it lets each hemisphere carry
its own amplitude, reference epoch, and shape parameters (a point we return to in
Section~\ref{sec:discussion}, because the hemispheres are related rather than
independent). For every measurement we derive an absolute latitude
$|\mathrm{lat}|$ and a decimal year, and for the amplitude proxy we sum the daily
corrected area of all groups above a $50~\msh$ noise floor within each
hemisphere and smooth the result with a rolling window.

% ============================================================================
\section{The classical universal-wing model}
\label{sec:classical}
% ============================================================================

A single butterfly wing, viewed hemisphere by hemisphere, is a distribution of
absolute emergence latitudes that drifts equatorward and narrows as the cycle
ages. Our aim in this section is to reduce that wing to as few numbers as the
data will allow, separating what is universal across cycles from what is
specific to a given cycle, and to do so in a form that can later serve as the
baseline density for the diffusion model of Section~\ref{sec:diffusion}.

\subsection{Time alignment and the universal mean path}
\label{sec:meanpath}

The natural clock for a wing is not the calendar but the cycle's own phase, and
so we align every hemicycle to a common reference epoch \tzero{} defined as the
decimal year at which the yearly-mean absolute latitude crosses $15\degr$ on its
equatorward descent, obtained by linear interpolation between the last year above
and the first year below (a physically meaningful anchor, since it sits on the
well-aligned decay phase rather than at the poorly-defined cycle onset).
Measuring standardized time as $\tau =
\mathrm{decimal\ year} - \tzero$, we find that the mean latitude of every
hemicycle collapses onto a single exponential path,
%
\begin{equation}
  \mu(\tau) = a\,e^{-\tau/b},
  \label{eq:meanpath}
\end{equation}
%
which we fit once by pooling the binned $(\tau,\mu)$ points from all hemicycles.
The universal path has an amplitude of $a = 16.31\degr$ (95\% CI
$[16.25, 16.44]$) at $\tau=0$ and an e-folding time of $b = 8.14$~yr
$[8.01, 8.25]$; these two numbers, together, are the universal skeleton of the
mean drift.

The $15\degr$ crossing is a convenient but imperfect anchor, and so we refine
\tzero{} per hemicycle in two steps: a first alignment shift $\delta_{S1}$ that
minimizes each hemicycle's residual against the universal template, and a second
shift $\Delta t_{S3}$ (bounded to $\pm2$~yr) obtained during the likelihood
optimization described in Section~\ref{sec:fitting}. The single quantity our
downstream code consumes is the effective reference epoch $t_0^{\rm total} =
t_0^{15\degr} + \delta_{S1} + \Delta t_{S3}$; for a hemicycle we have never seen,
$\Delta t_{S3}$ is by construction unidentifiable, and we fall back on the raw
$15\degr$ crossing as a proxy.

\subsection{The width model}
\label{sec:width}

The width of the wing is where the physics of the emergence zone lives, and it
is not monotonic in either time or latitude: few spots emerge at high latitude
early in the cycle, the zone is broadest at mid-cycle, and it collapses to a thin
band near the equator at the end. Because $\mu(\tau)$ is monotonic and therefore
invertible, we found it more physical to model the spread as a function of
\emph{where} the wing sits rather than \emph{when}, remapping the width to
$\sigma(\mu)$ (that is, the spread at a given mean latitude, independent of the
cycle's clock). In this coordinate the width is well described by two straight
lines joined at a peak latitude \mupeak,
%
\begin{equation}
  \sigma(\mu) =
  \begin{cases}
    m_{\rm shared}\,\mu + b_{\rm shared}, & \mu \leq \mupeak \\[4pt]
    m_i\,\mu + b_{\rm pol}, & \mu > \mupeak,
  \end{cases}
  \label{eq:sigma}
\end{equation}
%
where continuity at \mupeak{} fixes $b_{\rm pol}$ and the result is clipped to be
non-negative.

The equatorward arm ($\mu \leq \mupeak$) is the headline result of this section:
its slope and intercept, $m_{\rm shared} = 0.452$~$\degr\,\degr^{-1}$
$[0.438, 0.460]$ and $b_{\rm shared} = 0.16\degr$ $[0.08, 0.30]$, are
\emph{shared by every hemicycle}, so that the equatorward collapse of the active
zone follows one and the same line regardless of cycle strength. We interpret
this shared line as evidence that the equatorward narrowing is a
dynamo-controlled process common to all cycles, in the same family of universal
behaviors as the standard drift path of \citet[CIT]{Hathaway2011} and the
universal decline of \citet[CIT]{CameronSchussler2016}. The poleward arm ($\mu >
\mupeak$), in contrast, is cycle-specific: its slope $m_i$ and the detachment
latitude \mupeak{} both scale with cycle amplitude, so that stronger cycles
detach at higher latitude and descend more steeply.

\subsection{Amplitude coupling}
\label{sec:amplitude}

The bridge from a descriptive model to a generative one is the coupling between
cycle \emph{strength} and wing \emph{shape}, and we make that coupling explicit
by regressing the three cycle-specific parameters on the hemispheric amplitude
$A$ (the smoothed peak corrected area, in $\msh/1000$). All three relationships
are consistent with a simple linear form $\mathrm{param}(A) = a\,A + b$: the
detachment latitude $\mupeak(A)$ has slope $a=26.75$ and intercept $b=10.27\degr$;
the poleward width slope $m_i(A)$ has slope $a=-0.885$ and intercept $b=-0.034$;
and the high-latitude activity gate $\muzero(A)$ (the latitude above which we
treat the wing as not yet active, so that a year is included only when
$\mu(\tau) \leq \muzero(A)$) has slope $a=17.18$ and intercept $b=21.83\degr$. A
year is scored under the model only when it is active.

Two contrasts with SARG are worth drawing here. SARG fixes the emergence centroid
at $28\degr$ and lets only the band width vary, and it ties that width to the
whole-Sun sunspot number \citep[CIT]{Jha2025}; our model, by comparison, lets the
poleward detachment latitude and the poleward width both rise with amplitude, and
it ties them to the \emph{hemispheric} area rather than a global index. As a
direct check of the SARG width prescription against our directly measured areas,
we fit the single free constant $k$ in the SARG-style law $\sigma_\lambda(A) =
1.5\degr + 3.8\degr\,[1 - e^{-A/k}]$ and obtain $k \approx 8.2~\msh$, far below
the value of $400$ used in the SARG papers; this difference is expected, because
our catalog measures corrected area directly (with a median of order $100~\msh$),
whereas SARG derives its areas from the sunspot number through a fixed
conversion. % TODO(andres): confirm you want to keep the direct k comparison in
% the main text vs. moving it to an appendix.

\subsection{Fitting and validation}
\label{sec:fitting}

The full model has eight global parameters (the two mean-path parameters, the
two equatorward-arm parameters, and the four amplitude-regression coefficients
for \mupeak{} and $m_i$), plus the per-hemicycle timeshifts, and we fit them by
maximizing the per-block Gaussian log-likelihood in three stages. The first two
stages use Nelder-Mead (rather than a gradient method) because the activity gate
$\mu(\tau) \leq \muzero(A)$ is discontinuous in the parameters that move the mean
path, and the third stage uses L-BFGS-B on all globals plus the per-hemicycle
$\Delta t_{S3}$, warm-started from the second stage and made differentiable by
replacing the hard gate with a soft sigmoid ($T = 0.5\degr$) during optimization
while scoring with the hard gate. We add a coverage penalty to the objective so
that the optimizer cannot buy likelihood by quietly shrinking the fraction of
observed spots the model actually covers.

Our scoreboard metric is the mean per-year-normalized negative log-likelihood
(NLL), in nats per year-block, which weights each phase of the cycle equally
rather than letting the spot-rich maximum years dominate. On all $26$ hemicycles
the training NLL is $3.071$ nats per year-block $[3.063, 3.077]$ at a coverage of
$0.963$ $[0.961, 0.973]$. To ask how much of this fit is genuine rather than
overfit, we ran leave-one-cycle-out (LOCO) cross-validation, refitting the entire
pipeline on $25$ hemicycles and scoring the held-out one under its raw $15\degr$
epoch (with no test-time timeshift, which would leak). We also tested three
structural extensions of the model class against this baseline: a Student-$t$
likelihood with an asymmetric width (test NLL $+0.007$), a Student-$t$ likelihood
alone (test NLL $0.000$), and power-law-with-offset amplitude regressions (test
NLL $+0.016$). None of the three improved the held-out likelihood, and we
therefore interpret $3.087$ nats per year-block as the generalization floor of
this model class. We quantify parameter uncertainty with an $m$-out-of-$n$
bootstrap ($m = 0.95n$ hemicycles, $200$ draws, refitting the full pipeline for
each draw), and it is these draws that yield the confidence intervals quoted
above and in Table~\ref{tab:params}.

% TODO(andres): the LOCO ``Variant A'' naming and the coverage-penalty design
% are yours; confirm the framing of 3.087 as a ``generalization floor'' matches
% how you want to present it.

\begin{deluxetable}{lcc}
\tablecaption{Classical model parameters (all-data fit, $N=26$ hemicycles) with
bootstrap 95\% confidence intervals.\label{tab:params}}
\tablehead{\colhead{Parameter} & \colhead{Point} & \colhead{95\% CI}}
\startdata
$a$ (mean path, $\mu(\tau{=}0)$) [\degr]      & $16.31$   & $[16.25, 16.44]$ \\
$b$ (mean-path e-folding) [yr]                & $8.14$    & $[8.01, 8.25]$   \\
$m_{\rm shared}$ (eq.\ arm slope)             & $0.452$   & $[0.438, 0.460]$ \\
$b_{\rm shared}$ (eq.\ arm intercept) [\degr] & $0.16$    & $[0.08, 0.30]$   \\
$a_{\mupeak}$ (\mupeak\ vs $A$)               & $26.75$   & $[24.08, 28.57]$ \\
$b_{\mupeak}$ (\mupeak\ at $A{=}0$) [\degr]   & $10.27$   & $[9.97, 10.80]$  \\
$a_{m_i}$ ($m_i$ vs $A$)                      & $-0.885$  & $[-1.168, -0.617]$ \\
$b_{m_i}$ ($m_i$ at $A{=}0$)                  & $-0.034$  & $[-0.089, 0.022]$ \\
$a_{\muzero}$ (\muzero\ vs $A$)               & $17.18$   & $[14.62, 19.38]$ \\
$b_{\muzero}$ (\muzero\ at $A{=}0$) [\degr]   & $21.83$   & $[21.37, 22.42]$ \\
NLL (training) [nats/yr-block]               & $3.071$   & $[3.063, 3.077]$ \\
Coverage (training)                          & $0.963$   & $[0.961, 0.973]$ \\
\enddata
\end{deluxetable}

% ============================================================================
\section{Score-based residual refinement}
\label{sec:diffusion}
% ============================================================================

The classical model is smooth by construction, and a real wing is not: at any
given phase the observed latitude density carries structure (asymmetry, mild
bimodality, phase-dependent skew) that a single Gaussian cannot represent. Our
second modeling step is to learn that leftover structure with a score-based
diffusion model, keeping the classical Gaussian as a physically interpretable
baseline and letting the diffusion model supply only the correction.

\subsection{Target and setup}
\label{sec:target}

We slice each hemicycle into non-overlapping six-month windows (discarding any
window with fewer than $20$ emergences), and in each window we represent the
latitude distribution as a density on a fixed grid of $15$ bins of $3\degr$ each,
spanning $|\mathrm{lat}|$ from $0\degr$ to $45\degr$. The diffusion target is the
\emph{residual} density, defined as the empirical density minus the classical
parametric density (the latter obtained by integrating the model Gaussian over
each bin). We split the data by hemicycle rather than by window, because the
windows of one hemicycle share an amplitude and a shape and a row-wise split
would leak that information across the train/test boundary; sorting the
hemicycles by amplitude and taking every fifth one gives us a test set of six
hemicycles (12N, 13S, 15S, 19N, 22S, 24S) that we reserve entirely, fitting
nothing on them until the final comparison of Section~\ref{sec:test}.

\subsection{Model and training}
\label{sec:model}

Because the target is a $15$-dimensional vector rather than an image, the
score network is a multilayer perceptron (three hidden layers of $128$ units with
SiLU activations) rather than a convolutional U-Net, taking the noised residual
and a sinusoidal timestep embedding as input and predicting the added noise (the
standard $\epsilon$-prediction parameterization). We use a variance-preserving
cosine noise schedule \citep[$T=200$;][CIT]{NicholDhariwal2021}
and train with the standard $\epsilon$-prediction objective (mean-squared error
between the predicted and the true noise), using AdamW at a learning rate of
$10^{-3}$. Samples are drawn with a deterministic DDIM sampler and added to the
classical density to form the combined prediction.

The conditioning is where the physics enters the network. Our base conditioning
vector is four-dimensional, comprising the smoothed area, the universal mean
latitude $\mu(\tau)$, the model width $\sigma$, and the cycle amplitude; we also
tried richer conditioning sets (cycle number and hemisphere identity, the
contemporaneous opposite-hemisphere amplitude and phase, and an autoregressive
history of the hemicycle's own area), two conditioning mechanisms (raw
concatenation and feature-wise linear modulation, FiLM), an optional Fourier
feature lifting of the conditioning scalars, and classifier-free guidance. One
choice deserves emphasis: we select the final checkpoint on a distributional
metric (the validation Earth Mover's distance) rather than on the denoising
loss, because we found the denoising loss to rise even as the generated samples
became visibly more realistic.

\subsection{Metrics}
\label{sec:metrics}

We evaluate the model along two axes that, as it turns out, disagree. The first
is the hard-gated NLL, scored exactly as in the classical model (each observed
latitude scored under the combined density where the wing is active). We compute
it in two forms: a raw form that scores only the occupied bins, and a normalized
form that adds the log-partition term $\log Z$ so that mass spent outside the
occupied bins is properly charged; the distinction between the two turns out to
matter a great deal (Section~\ref{sec:guidance}). The second axis is a suite of
distributional metrics computed between the generated and the observed densities:
the Earth Mover's distance (a Wasserstein-1 distance, in degrees of latitude),
an energy distance, the continuous ranked probability score (CRPS), the error in
the recovered spread ($\sigma$-MAE), and a Talagrand rank histogram for
calibration. Finally, we use an oracle diagnostic (an auxiliary network that maps
the conditioning vector directly to a per-bin Gaussian) to estimate the best
NLL any model could achieve from a given conditioning set, which gives us an
information upper bound to compare the diffusion model against.

% ============================================================================
\section{Results}
\label{sec:results}
% ============================================================================

\subsection{The ablation, and a tension between metrics}
\label{sec:ablation}

Table~\ref{tab:val} shows the validation-set comparison across the conditioning
and architecture variants, and the first thing it reveals is that no diffusion
variant beats the classical model on NLL: the best NLL among the diffusion
models ($3.402$) is still above the classical baseline ($3.240$). However, the
same table shows that several variants beat the classical model handily on
distributional shape. The simplest conditioned model (E0: the four-dimensional
base conditioning with raw concatenation, no guidance, no Fourier features)
lowers the Earth Mover's distance to $1.365\degr$ from the classical $1.575\degr$,
lowers the CRPS of the mean latitude to $0.794$ from $0.963$, and roughly halves
the spread error ($\sigma$-MAE $0.622$ versus $1.153$). The added complexity we
tried (FiLM, Fourier features, extra conditioning groups) either failed to help
or actively hurt these distributional metrics, even when it lowered the raw NLL,
which is the first of several signs that NLL and distributional fidelity are
ranking the models differently.

\begin{deluxetable}{lccc}
\tablecaption{Validation-set comparison (55 blocks; $K{=}100$ samples per block).
Lower is better throughout. E0 is the deployed hybrid variant.\label{tab:val}}
\tablehead{\colhead{Model} & \colhead{NLL} & \colhead{EMD [\degr]} &
\colhead{$\sigma$-MAE [\degr]}}
\startdata
Classical (baseline)          & $\mathbf{3.240}$ & $1.575$          & $1.153$          \\
E0 (base, concat)             & $3.502$          & $\mathbf{1.365}$ & $\mathbf{0.622}$ \\
E1 (base + cycle/hemi)        & $3.457$          & $1.437$          & $0.630$          \\
E3 (base + trajectory)        & $3.437$          & $1.554$          & $0.696$          \\
E13 (base, FiLM + Fourier)    & $3.402$          & $1.944$          & $0.828$          \\
E5 (base + opp., FiLM)        & $4.014$          & $2.336$          & $0.746$          \\
\enddata
\tablecomments{Full 18-experiment sweep in the supplementary material.}
% TODO(andres): decide whether all 18 experiments go in an appendix table.
\end{deluxetable}

\subsection{Classifier-free guidance fails, and why}
\label{sec:guidance}

Classifier-free guidance is the standard lever for sharpening conditional
diffusion samples, and we expected it to help --- it did not, on any axis. Across
every guided variant the sweep verdict was the same (the unguided setting,
$w=0$, was best), and raising the guidance weight monotonically \emph{worsened}
the Earth Mover's distance rather than improving it. The reason is instructive
and worth stating plainly, because it is a property of the residual framing
rather than a bug in the sampler. The raw hard-gated NLL scores the combined
density only at the occupied bins, so guidance can drive the raw NLL down simply
by piling residual mass onto the occupied bins, while the mass it over-spends in
the empty bins is floored at zero and never charged; the normalized NLL, which
adds $\log Z$, closes this leak, and under guidance the gap between the raw and
the normalized NLL widens exactly as this picture predicts. Our sampling
diagnostics confirm the mechanism directly: as the guidance weight rises, the
added mass climbs, the ensemble diversity either collapses toward mode collapse
or explodes into a handful of enormous samples, and a growing fraction of the
observed latitudes fall into bins the model has driven to negative density. We
take this as evidence that guidance buys nothing here and actively destabilizes
the ensemble, and we deploy the model unguided.

\subsection{An information ceiling}
\label{sec:ceiling}

The failure of every diffusion variant to beat the classical model on NLL raised
an obvious question: is the limitation in the architecture, in the residual
framing, or in the conditioning information itself? The oracle diagnostic answers
it. Across conditioning sets the oracle NLL sits between roughly $3.46$ and
$3.65$, which is close to both the classical baseline and the diffusion model's
own NLL, and the smallness of the diffusion-to-oracle gap tells us that the
diffusion model is already extracting nearly all of the residual signal the
conditioning set contains. In other words, the ceiling is set by the
\emph{information} in the covariates, not by the capacity of the network. To test
this interpretation we reframed the target, replacing the residual with the
empirical density itself modeled in logit space (so that non-negativity and the
unit-mass constraint hold by construction), with the classical density demoted
from a fixed pedestal to a Kullback-Leibler prior. The reframing narrowed the gap
to the classical NLL (the best empirical-target variant reaches $3.373$ versus
the best residual-target $3.402$). However, it did not break through it, and the
empirical model sits essentially at its own oracle (its NLL and its oracle NLL
agree to within the sampling scatter). We interpret the two results
together as evidence that the residual signal available in the present
conditioning set is nearly exhausted, and that beating the classical model on
likelihood will require genuinely new predictive covariates rather than a better
network or a different target.

\subsection{The held-out test comparison}
\label{sec:test}

The real test of a generative model is on cycles it has never seen, and so we
turn to the six reserved hemicycles, comparing four models on the pooled
held-out windows: the empirical distribution, an independent SARG-style baseline
(with its own universal path and amplitude-width law fit on the training split
only), the classical model, and the hybrid ``classical $+$ AI'' model (the E0
residual added to the classical density). Table~\ref{tab:test} gives the
verdict. On genuinely held-out cycles the hybrid model delivers substantially
better distributional fidelity than both the classical model and the SARG
baseline: the Earth Mover's distance improves by $24\%$ over the classical model
($1.391\degr$ versus $1.832\degr$), the energy distance by $28\%$, the CRPS of
the mean by $29\%$, and the spread error is roughly halved ($0.714\degr$ versus
$1.439\degr$). The cost, as on the validation set, is the pointwise NLL, which is
worse for the hybrid model ($3.588$) than for either baseline ($3.234$ and
$3.248$); this is a trade we make deliberately (and one we defend at length in
Section~\ref{sec:discussion}).

\begin{deluxetable}{lcccc}
\tablecaption{Held-out test comparison on the six reserved hemicycles (86
windows). Lower is better throughout.\label{tab:test}}
\tablehead{\colhead{Model} & \colhead{NLL} & \colhead{EMD [\degr]} &
\colhead{Energy} & \colhead{$\sigma$-MAE [\degr]}}
\startdata
SARG (independent baseline)   & $3.248$          & $1.531$          & $0.104$          & $1.114$          \\
Classical                     & $\mathbf{3.234}$ & $1.832$          & $0.106$          & $1.439$          \\
Classical $+$ AI (E0)         & $3.588$          & $\mathbf{1.391}$ & $\mathbf{0.077}$ & $\mathbf{0.714}$ \\
\enddata
\end{deluxetable}

\subsection{Physical plausibility}
\label{sec:physical}

A generative model of the butterfly diagram is only useful if its samples obey
the physics, and so we check each variant against Sp\"orer's law and the bounds
of the emergence zone. Every variant reproduces the equatorward drift with the
correct sign, and between $88\%$ and $93\%$ of the generated mass falls inside
the $5\degr$--$40\degr$ Sp\"orer band. However, the diffusion samples share one
systematic imperfection, in that they drift slightly too slowly: the observed
centroid slope is
$-2.19\degr\,\mathrm{yr}^{-1}$, the classical model gives $-1.72\degr\,
\mathrm{yr}^{-1}$, and the E0 hybrid gives $-1.52\degr\,\mathrm{yr}^{-1}$, so the
model under-predicts the drift rate while preserving its direction. We flag this
as the clearest target for future improvement.

% ============================================================================
\section{Discussion}
\label{sec:discussion}
% ============================================================================

The central result of this work is also its central tension: the hybrid model
generates more realistic butterfly wings than the classical model, and yet it
scores worse on likelihood. This is not a contradiction so much as a statement
about what the two metrics reward. The NLL is dominated by the tails of the
density and is minimized by a smooth, slightly over-dispersed Gaussian that never
assigns a very low probability to any observed spot; the distributional metrics
(the Earth Mover's distance, the energy distance, the CRPS, and the spread error)
instead reward getting the \emph{shape} of the density right, which is precisely
what the diffusion model adds. For a \emph{generative} model, whose purpose is to
draw new samples that look like real cycles rather than to assign the highest
possible probability to the observed ones, we argue that the distributional
metrics are the more appropriate arbiter, and that the likelihood penalty is a
price worth paying for shape realism. We do not claim this settles the question,
and we suspect the right long-term answer is a proper scoring rule that credits
shape and calibration together.

The residual framing carries an implicit assumption that is worth naming: it
assumes that the conditioning covariates carry information about the residual
\emph{beyond} what the classical model has already extracted from those same
covariates. However, our oracle diagnostic shows that this assumption is only
weakly satisfied with the present conditioning set, because the oracle NLL sits
close to
both the classical baseline and the diffusion model's own likelihood. We
interpret the failure of classifier-free guidance in the same light: guidance
sharpens a conditional distribution around information the conditioning already
carries, and where that information is nearly exhausted there is nothing left for
guidance to sharpen, so it can only exploit the scoring leak or destabilize the
ensemble. Rather than read these as failures of the method, we read them as a
measurement of how much predictive signal the standard cycle parameters actually
contain about the fine structure of a wing, and the answer is: less than one
might hope.

Positioned against the existing generators, our model makes two moves that we
believe are new for the butterfly diagram. The first is that we tie the full
shape of each wing to that hemisphere's own amplitude, rather than fixing the
centroid and tying only the width to a global index as SARG does, or using a
whole-cycle amplitude as \citet[CIT]{Jiang2011} do; the observed hemispheric
asymmetry is bounded (the hemispheres are roughly $80\%$ coupled;
\citealt[CIT]{Norton2014}), which is why we treat them as related processes with
independent amplitudes rather than as fully independent ones.
% TODO(andres): confirm you want to lean on the Norton 80%-coupling number here.
The second is the residual score-based refinement of a physical baseline, which
is well established in atmospheric downscaling \citep[CIT]{Mardani2025} but which,
to the best of our knowledge, has not previously been applied in solar physics.
The clearest paths to improvement follow directly from the information ceiling we
found: new covariates (for example, the polar-field or terminator timing that
sets a cycle's onset, or the opposite hemisphere's fine structure) rather than a
larger network, and a drift-rate-aware objective to cure the shallow-drift bias
of Section~\ref{sec:physical}.

% ============================================================================
\section{Summary}
\label{sec:summary}
% ============================================================================

We have built a two-part generative model of the solar butterfly diagram. The
first part is a compact classical description of a butterfly wing, in which the
mean emergence latitude follows a universal exponential path
($a=16.31\degr$, $b=8.14$~yr), the equatorward collapse of the width follows a
single line shared by every cycle, and the poleward shape and the activity gate
scale with each hemisphere's amplitude; leave-one-cycle-out validation places the
generalization floor of this model at $3.087$ nats per year-block. The second
part is a score-based diffusion model that generates the residual density the
classical Gaussian leaves behind, conditioned on cycle amplitude and phase. On
six hemispheric cycles held out from all fitting, the hybrid model improves the
Earth Mover's distance by $24\%$, the energy distance by $28\%$, and the spread
error by roughly a factor of two over both the classical model and an independent
SARG-style baseline, at the cost of a worse pointwise likelihood that continues
to favor the smoother Gaussian. We find, through an oracle diagnostic, that this
likelihood ceiling is set by the information in the conditioning covariates rather
than by the network or the target framing, and we find that classifier-free
guidance offers no help and destabilizes the ensemble. To the best of our
knowledge, this is the first score-based generative model of the solar butterfly
diagram, and it establishes both a working hybrid architecture and a clear
measurement of where its predictive information runs out.

% ============================================================================
\section*{Acknowledgements}
% ============================================================================

This work was carried out within an eleven-week undergraduate research program at
the COFFIES Science Center (approximately $40$ students), which developed the
classical model, the diffusion pipeline, and the evaluation framework described
here. % TODO(andres): full acknowledgement wording, funding, and data
% provenance (survey list and DOIs for the composite catalog).

\facility{}

% \bibliography{butterflai}   % TODO(andres): supply the .bib

\end{document}
